% EGA V, English translation by P. Blass, 1987 (cote 24): archive cote 24, site page 59 (PDF page 60).
% typescript. Draft compared: none.
% AI-transcribed by Claude Opus 5.5 (claude-opus-5-5, adaptive thinking, effort high), Message Batches API, batch ega24b (msgbatch_01RGJJeHdR96U16ZWFuUjtDx), 2026-11-13; not yet checked by eye.
\scanpage{24/59}
\mspage{59}

that $Y_{\xi}$ should be geometrically irreducible or geometrically
integral, and the two problems [(\uline{respé et non respé}) Fr p 50]
considered in 8.18 are therefore equivalent. On the other hand,
by 5.6, the set U of $\xi \in P$ such that $Y_{\xi}$ is geometrically integral
of dim d-1 (the dimension hypothesis stating that $\phi_{\xi}$ is $X_{k(\xi)}$
regular) is \uline{open}. We will exhibit a non-empty open \uline{evident} subset
P-Z contained in U and taking for Z the \uline{union of} $g(\underline{Y}^{1\ \mathrm{sing}}$ and of
the hyperplanes $H_{x}$ of $P^{v}$ defined by the f(x), $X \in T$. For $\xi \in P^{v}$ -
Z, $Y_{\xi}$ is smooth of dimension
(d - 1) and since it is geometrically connected by 6.1 it is
geometrically integral. We have to, therefore, express (explain)
(prove) that every irreducible component of codimension one of Z
meets the open set U. But these irreducible components are the $H_{x}$
[they are repeated possibly, but it is not essential] and also $g(Y^{1}$
$^{\mathrm{sing}})$ when the latter are indeed of codimension one i.e. $X^{1}$ ``not
ruled'' for f (Nota Bene: we use the irreducibility of $Y^{1\mathrm{sing}}$). On
the other hand, in order that this latter set should meet the open
set U it is necessary and sufficient that $g(Y^{1\mathrm{sing}})$ which contains
an open and dense set) should meet U. This proves 8.18 in this
case. If we do not suppose that X is normal, we apply the
previous result to the normalization of X the reasoning is
immediate and I do not give the details here. N.B. In the \uline{case}
[respé] 8.18 is contained in 8.16 more precisely except in the
case d = 2. It is for the case \uline{[non-respé] that it may be better}
\uline{not to require]} d = 2 and not only $d \geq 2$ .....

It remains to explain, (make explicit) the conditions a) and
b) of 8.18. This leads us to exmaine in a general way the
